% Chapter 1: Introduction & Overview — Mid-Term Report Phase 3

\chapter{Introduction and Overview}

\section{Problem Statement and Motivation}

Academic credentials (e.g., diplomas, degrees, certificates, and academic transcripts) serve as fundamental attestations of an individual's knowledge, skills, and qualifications. In modern hiring and education pipelines, these documents represent critical evidence for employers and academic institutions during candidate evaluations. However, the legacy certificate issuance and verification ecosystem suffers from severe structural vulnerabilities:

\begin{itemize}
    \item \textbf{Credential Forgery and Fraud}: Physical paper degrees and unauthenticated scanned PDFs can be easily duplicated, manipulated, or fabricated using modern digital editing and printing tools.
    \item \textbf{High Verification Overhead and Lack of Direct Verifiability}: Cross-institution verification often requires manual inquiries, physical stamp verifications, or dependence on centralized registrar portals, which are slow, costly, and subject to single-point-of-failure bottlenecks.
    \item \textbf{Over-disclosure of Sensitive Information}: Traditional certificates mandate revealing the entire document (including Personally Identifiable Information (PII) such as date of birth, student ID, detailed GPA) even when a verifier only needs to assert whether the candidate graduated.
\end{itemize}

The \textbf{BK Credential System} is engineered as a three-tier hybrid architecture combining a relational database, an on-chain blockchain trust anchor, and decentralized storage (IPFS), enabling fully trustless verification without relying on centralized backend servers. The system addresses these challenges by guaranteeing:
\begin{itemize}
    \item \textbf{Tamper-evidence}: Enforced through cryptographic hashing and public blockchain state anchoring.
    \item \textbf{Privacy-Preserving Selective Disclosure}: Permitting credential holders to disclose only required attributes without compromising signature validity.
    \item \textbf{Cost and Scalability Efficiency}: Utilizing Merkle tree batching to amortize blockchain transaction costs across thousands of graduating students.
    \item \textbf{Trustless Verification}: Enabling any third party to cryptographically verify credentials directly against blockchain trust anchors without depending on proprietary university backend infrastructure.
\end{itemize}

In Phase 2 (HK252-DATN-131), we designed and implemented a functional prototype leveraging \textit{SD-JWT VC (Selective Disclosure JSON Web Token Verifiable Credential)} paired with an on-chain \textit{Merkle Tree} registry. However, rigorous evaluation at university scale ($\geq$10,000 graduates per cohort) revealed two significant operational bottlenecks:

\begin{enumerate}
    \item \textbf{IPFS Operational Overhead}: Mandatory pinning of Public Canonical Views (PCV) per credential to IPFS incurred recurring storage fees (\$0.001--\$0.01/credential/month) and introduced an operational dependency on external IPFS gateways.
    \item \textbf{SD-JWT Library Complexity}: Lack of native Web3 standard support for the ES256K curve in standard JWT libraries necessitated custom SD-JWT implementations with substantial JWT envelope and \texttt{\_sd\_alg} parsing overhead.
\end{enumerate}

\section{Phase 3 Objectives}

Phase 3 transitions the architecture to \textbf{EIP-712 (Typed Structured Data Hashing and Signing)}, retaining all cryptographic guarantees of Phase 2 while resolving operational friction. The explicit project objectives are detailed in Table~\ref{tab:objectives}.

\begin{table}[H]
\centering
\renewcommand{\arraystretch}{1.3}
\begin{tabular}{|c|p{5cm}|p{8.5cm}|}
\hline
\textbf{\#} & \textbf{Objective} & \textbf{Technical Approach} \\
\hline
O1 & Eliminate mandatory IPFS dependencies & Transition entire credential payload to structured off-chain JSON with EIP-712 typed signature. \\
\hline
O2 & Simplify cryptographic signing/verification & Leverage standard, battle-tested \texttt{ethers.signTypedData} / \texttt{ethers.verifyTypedData} primitives. \\
\hline
O3 & Reduce on-chain gas overhead by $\geq$10$\times$ & Perform EIP-712 signing fully off-chain (zero gas per individual credential) with a single \texttt{anchorBatch} transaction per cohort. \\
\hline
O4 & Streamlined Selective Disclosure & Integrate salted-hash keccak256 commitments directly into the EIP-712 struct payload. \\
\hline
O5 & Gas-optimized Revocation Management & Deploy on-chain Bitmap revocation (1 bit/credential, saving $\sim$64$\times$ storage relative to per-credential mappings). \\
\hline
O6 & High-throughput Processing ($\geq$500 cred/sec) & Develop a lightweight, modular TypeScript SDK optimized for bulk processing. \\
\hline
\end{tabular}
\caption{Phase 3 Key Project Objectives}
\label{tab:objectives}
\end{table}

\section{Scope of Phase 3}

Rather than serving as a linear incremental update, Phase 3 represents a fundamental architectural paradigm shift. It is formulated to validate the core design thesis: \textit{at institutional university scale, a decoupled architecture where the blockchain functions strictly as a Trust Anchor and cryptographic signatures are executed off-chain via EIP-712 achieves equivalent security and privacy guarantees while drastically minimizing operational costs and operational complexity}.

The technical scope of Phase 3 is structured across five systematic engineering layers:

\begin{enumerate}
    \item \textbf{Cryptographic and Protocol Layer}:
    \begin{itemize}
        \item Formal specification of the EIP-712 typed structured data schema, including domain separation parameters (\texttt{name}, \texttt{version}, \texttt{chainId}, \texttt{verifyingContract}) and nested data structures (\texttt{PublicClaims}, \texttt{BkCredential}).
        \item Design of a 256-bit salted-hash commitment scheme for privacy-preserving Selective Disclosure, computing attribute commitments as $C_i = \text{keccak256}(\text{salt}_i \parallel \text{key}_i \parallel \text{value}_i)$ directly within the signed EIP-712 payload.
        \item Integration of OpenZeppelin \texttt{StandardMerkleTree} with double keccak256 leaf hashing ($\text{leaf} = \text{keccak256}(\text{keccak256}(\text{abi.encode}(\dots)))$) to prevent second-preimage collision attacks.
    \end{itemize}
    
    \item \textbf{Smart Contract and Trust Anchor Layer}:
    \begin{itemize}
        \item Implementation and verification of \texttt{IssuerRegistry.sol}, governing authorized signing addresses, decentralized identifiers (DIDs), and two-step ownership management via OpenZeppelin \texttt{Ownable2Step}.
        \item Implementation of \texttt{CredentialRegistryV3.sol}, functioning as the on-chain Trust Anchor by recording batch Merkle root commitments via \texttt{anchorBatch()} and managing credential revocations through an optimized 256-bit bitmap mapping (\texttt{uint256[]}), supported by \texttt{Pausable} emergency controls.
    \end{itemize}

    \item \textbf{Software Development Kit (SDK) Layer}:
    \begin{itemize}
        \item Engineering of a standalone, modular TypeScript SDK (\texttt{@bkcred/sdk}) comprising 6 specialized packages: \texttt{domain}, \texttt{signer}, \texttt{verifier}, \texttt{disclosure}, \texttt{merkle}, and barrel export \texttt{index}.
        \item Formulation of an offline-capable, deterministic \textbf{10-Step Verification Pipeline} combining local schema validation, cryptographic struct hashing, commitment checking, Merkle proof evaluation, and online smart contract assertions.
    \end{itemize}

    \item \textbf{Full-Stack Application and Portal Layer}:
    \begin{itemize}
        \item Construction of 4 dedicated web portals on Next.js 14 (App Router) and PostgreSQL (Prisma ORM): (1) \textit{Issuer Portal} for batch CSV parsing, off-chain bulk signing, and on-chain root anchoring; (2) \textit{Holder Portal} featuring passwordless email authentication (Privy) and an interactive Selective Disclosure customization interface; (3) \textit{Public Verifier Portal} supporting direct JSON file validation and deep-link routing (\texttt{/v3/verify?credId=\dots}); and (4) \textit{Administrative Portal} for signer key lifecycle management.
    \end{itemize}

    \item \textbf{Empirical Evaluation and Benchmarking Layer}:
    \begin{itemize}
        \item Development of automated testing harness and fixture generators capable of producing cohorts of 10,000 synthetic graduate credentials.
        \item Execution of multi-scenario benchmark suites evaluating offline signing throughput, end-to-end issuance latency, concurrent verification capacity, gas profiling, and static security analysis using the Slither framework.
    \end{itemize}
\end{enumerate}

This mid-term report captures our progress through the completion of \textbf{Week 4} (Specification, Smart Contracts, SDK, and Web Portals), representing 100\% of the scheduled mid-term milestones within the 7-week roadmap.

\section{Comparative Analysis with the Three-Tier Hybrid SD-JWT VC Architecture (Phase 2)}

To substantiate the motivation for Phase 3, it is critical to contrast its foundational design decisions with the Phase 2 prototype (HK252-DATN-131). 

In Phase 2, the system adopted an on-chain asset and token-centric model where individual credentials were instantiated as ERC-721 Non-Fungible Tokens (NFTs). While this demonstrated basic blockchain immutability, it suffered from severe scalability bottlenecks:
\begin{itemize}
    \item \textbf{Linear Gas Cost Scaling}: Every issued certificate required an individual on-chain \texttt{mintCertificate()} transaction ($\sim$103,000 gas per certificate), escalating to over 1.03 billion gas for a standard university graduating cohort of 10,000 students.
    \item \textbf{Mandatory IPFS Pinning Overhead}: Each credential mandated pinning a Public Canonical View (PCV) to IPFS, incurring continuous pinning subscription fees ($\sim$\$0.001--\$0.01/credential/month) and creating a fragile operational dependency on external gateway availability.
    \item \textbf{Non-Standard Cryptographic Stacks}: Employing SD-JWT with ECDSA over secp256k1 (ES256K) lacked native support in standard OAuth/JWT libraries, necessitating a custom parser with heavy envelope formatting and complex \texttt{\_sd\_alg} metadata.
    \item \textbf{Storage-Inefficient Revocation}: Tracking revocation via \texttt{mapping(bytes32 => bool)} consumed an entire 32-byte storage slot per revoked credential.
\end{itemize}

Phase 3 resolves these bottlenecks by transitioning from a token-centric paradigm to a \textbf{pure attestation and Trust Anchor paradigm}. Under Phase 3:
\begin{itemize}
    \item \textbf{Zero-Gas Issuance}: Individual credentials exist purely off-chain as structured JSON documents cryptographically signed via native EIP-712 typed structured data.
    \item \textbf{Amortized On-Chain Anchor}: Only a single \texttt{anchorBatch} transaction ($\sim$105,000 gas total) is executed per cohort, binding up to $2^{32}$ certificates to the blockchain for approximately 10.5 gas per graduate.
    \item \textbf{Elimination of Mandatory IPFS}: Raw credential files are distributed directly to holders, eliminating per-credential IPFS pinning costs and gateway latency bottlenecks.
    \item \textbf{Native Web3 Selective Disclosure}: Attribute commitments utilize EVM-native keccak256 salted hashes verified in a single computation without custom JWT dependencies.
    \item \textbf{Bitmap Revocation}: Revocation bits are packed into 256-bit words, saving $\sim$64$\times$ in on-chain storage.
\end{itemize}

Table~\ref{tab:phase-comparison} provides a comprehensive multi-dimensional comparison between Phase 2 and Phase 3.

\begin{table}[H]
\centering
\renewcommand{\arraystretch}{1.3}
\begin{tabular}{|p{3.8cm}|p{5.4cm}|p{5.4cm}|}
\hline
\textbf{Evaluation Metric} & \textbf{Phase 2 (SD-JWT VC + NFT)} & \textbf{Phase 3 (EIP-712 Trust Anchor)} \\
\hline
Credential Paradigm & Token-centric (ERC-721 NFT) & Pure cryptographic attestation \\
\hline
Signature Protocol & ES256K SD-JWT (Compact format) & EIP-712 Typed Structured Data \\
\hline
Issuance Gas Cost & $\sim$103,000 gas per credential & 0 gas (Off-chain typed signing) \\
\hline
Batch Anchor Overhead & $N$ individual on-chain transactions & 1 tx \texttt{anchorBatch()} ($\sim$105k gas total) \\
\hline
Storage Architecture & Mandatory per-credential IPFS pin & Self-contained JSON (IPFS optional) \\
\hline
Selective Disclosure & SD-JWT \texttt{\_sd[]} array + SHA-256 & Salted keccak256 commitment array \\
\hline
Revocation Storage & \texttt{mapping(bytes32 => bool)} (1 slot/cred) & \texttt{uint256[]} Bitmap (256 creds/slot) \\
\hline
Verification Dependencies & Custom SD-JWT parser + JWT verify & Standard \texttt{ethers.js} v6 \\
\hline
RPC Query Footprint & 3 \texttt{eth\_call} queries + IPFS fetch & 2 \texttt{eth\_call} queries (Batchable) \\
\hline
End-to-End Latency & $\sim$608 ms (Dominated by IPFS) & $\sim$204 ms ($\sim$3$\times$ performance gain) \\
\hline
Web3 Wallet Need & Mandatory for student (Receive NFT) & Not required for credential holders \\
\hline
Total Cost (10k Batch) & $\sim$1.03B gas + \$50/month IPFS & $\sim$105,000 gas total + \$0 IPFS \\
\hline
\end{tabular}
\caption{Comprehensive Architectural and Operational Comparison: Phase 2 vs. Phase 3}
\label{tab:phase-comparison}
\end{table}

\section{Report Organization}

The remainder of this report is structured as follows:
\begin{itemize}
    \item \textbf{Chapter 2 (Theoretical Foundations)}: Supplements core cryptographic concepts with EIP-712 typed hashing, salted-hash selective disclosure, on-chain bitmap revocation, and OpenZeppelin StandardMerkleTree.
    \item \textbf{Chapter 3 (System Design)}: Outlines overall system architecture, credential type schemas, smart contract interfaces, SDK specifications, workflow pipelines, and database modeling.
    \item \textbf{Chapter 4 (System Implementation)}: Documents the technical implementation of smart contracts, SDK modules, web portals, and API integration.
    \item \textbf{Chapter 5 (Testing and Evaluation)}: Details smart contract coverage, gas measurements, SDK performance benchmarks, and static security audit findings.
    \item \textbf{Chapter 6 (Conclusion and Future Work)}: Summarizes mid-term milestones, presents preliminary comparative metrics, and defines the roadmap for Weeks 5--7.
\end{itemize}
